%! AP Statistics — Unit 1, Lesson 1.5 — Guided Notes ANSWER KEY
\documentclass[10pt]{article}
\usepackage{apstats-article}
\usepackage{apstats-key}

%=============================================================================
\begin{document}

\pageheader{Unit 1, Lesson 1.5}{Guided Notes --- Answer Key}
\begin{tabularx}{\linewidth}{X X X}
Name:\;\ans{Key} & Date:\; & Period:\;
\end{tabularx}

\vspace{0.12in}

\begin{objectivebox}
\small By the end of this lesson, I will be able to\dots\\[4pt]
\ansline{read and interpret dotplots, stemplots, and histograms for quantitative data.}\\
\ansline{identify when each display is most appropriate based on dataset size and purpose.}\\
\ansline{describe the overall pattern of a distribution (shape, center, spread) from a graph.}
\end{objectivebox}

\vspace{0.1in}

\begin{vocabbox}
\small
\noindent\textbf{\textcolor{navy}{Dotplot:}}\\
\ans{A graph in which each data value is shown as a dot above its position on a
number line. Reveals every individual value; best for small datasets.}\\[5pt]

\noindent\textbf{\textcolor{navy}{Stemplot (stem-and-leaf plot):}}\\
\ans{A display that separates each data value into a stem (leading digit(s)) and a
leaf (trailing digit). Preserves actual data values and sorts them automatically.}\\[5pt]

\noindent\textbf{\textcolor{navy}{Histogram:}}\\
\ans{A graph that groups quantitative data into equal-width intervals (bins) and
displays frequency or relative frequency as bar height. Bars touch.}\\[5pt]

\noindent\textbf{\textcolor{navy}{Bin (class interval):}}\\
\ans{An equal-width grouping of values on the horizontal axis of a histogram.
e.g., $[20, 40)$ means all values from 20 up to (but not including) 40.}\\[5pt]

\noindent\textbf{\textcolor{navy}{Gap / Cluster:}}\\
\ans{A gap is a range of values with no observations. A cluster is a concentration
of many observations in a narrow range.}\\[5pt]
\end{vocabbox}

\vspace{0.1in}

\begin{hookbox}
\small\textbf{Three displays, one dataset.}\\[6pt]
\ans{The dotplot and stemplot make it easy to see every individual value (e.g., that
five students are exactly 10 years old) and to identify the most common value
quickly. The histogram groups data and shows the overall shape --- useful for
seeing that most campers are 10--11 years old --- but hides the fact that the
mode is exactly 10. Each trade-off matters depending on the question.}
\end{hookbox}

\vspace{0.1in}

\begin{notesbox}{Display 1 --- Dotplots (Key)}
\small
\begin{multicols}{2}

\textbf{How to read a dotplot:}
\begin{itemize}[itemsep=3pt]
  \item Horizontal axis: \ans{values / scale} of the variable.
  \item Each dot: \ans{one individual (data point)}.
  \item Stack height: \ans{frequency (count)} at that value.
  \item Look for: center, spread, shape, \ans{gaps}, \ans{clusters}.
\end{itemize}

\columnbreak

\textbf{Best for:}
\begin{itemize}[itemsep=3pt]
  \item \ans{Small} datasets (roughly $n \le$ \ans{30--50})
  \item When you want to see \ans{all individual} data values
  \item Identifying gaps, clusters, and outliers quickly
\end{itemize}

\vspace{6pt}
\textbf{Limitation:}\\
Impractical when $n$ is \ans{large} --- dots pile up and overlap.

\end{multicols}
\end{notesbox}

\vspace{0.1in}

\begin{notesbox}{Display 2 --- Stemplots (Key)}
\small
\begin{multicols}{2}

\textbf{Structure:}
\begin{itemize}[itemsep=3pt]
  \item Stem = \ans{leading (tens or higher)} digits
  \item Leaf = \ans{final (units)} digit
  \item Leaves in \ans{ascending} order.
  \item Always include a \ans{key}: e.g., \textbf{7 $|$ 4} = 74.
\end{itemize}

\vspace{6pt}
\textbf{Completed stemplot} for 54, 58, 61, 65, 67, 72, 78, 79, 85:

\vspace{4pt}
\begin{tabular}{r|l}
\textbf{Stem} & \textbf{Leaf} \\
\hline
5 & \ans{4\; 8} \\
6 & \ans{1\; 5\; 7} \\
7 & \ans{2\; 8\; 9} \\
8 & \ans{5} \\
\end{tabular}

\columnbreak

\textbf{Best for:}
\begin{itemize}[itemsep=3pt]
  \item Moderately sized datasets ($n \approx$ \ans{15--50})
  \item When you want to \ans{preserve / read back} the actual data values
  \item Comparing two groups with a \ans{back-to-back} stemplot
\end{itemize}

\vspace{6pt}
\textbf{Advantage over dotplot:}\\
Data is automatically \ans{sorted}, making it easy to find the median.

\vspace{6pt}
\textbf{Limitation:}\\
Difficult with very large data values or non-integer data.

\end{multicols}
\end{notesbox}

\newpage

\begin{notesbox}{Display 3 --- Histograms (Key)}
\small
\begin{multicols}{2}

\textbf{Key features:}
\begin{itemize}[itemsep=3pt]
  \item Data grouped into equal-width \ans{bins / intervals}.
  \item Bar height = \ans{frequency} (or relative frequency).
  \item Bars \ans{touch} --- no gap between adjacent bins.
  \item Horizontal axis: the \ans{variable} being measured.
  \item Vertical axis must start at \ans{0}.
\end{itemize}

\vspace{6pt}
\textbf{Reading a histogram:}\\
Shows the \ans{shape / overall pattern} --- but NOT individual values.

\columnbreak

\textbf{Histogram vs.\ bar chart:}

\vspace{4pt}
\renewcommand{\arraystretch}{1.5}
\begin{tabularx}{0.46\textwidth}{@{} >{\bfseries}p{0.18\textwidth} X @{}}
\rowcolor{sky}
 & \textbf{Key difference} \\
\hline
Bars touch & \ans{quantitative} data \\
Bars separate & \ans{categorical} data \\
Order matters & \ans{histogram/dotplot} \\
Order arbitrary & \ans{bar chart} \\
\end{tabularx}

\vspace{6pt}
\textbf{Best for:}\\
Large datasets ($n > $ \ans{50}) where individual values matter less than overall \ans{shape}.

\end{multicols}
\end{notesbox}

\vspace{0.1in}

\begin{notesbox}{Worked Example --- Key}
\small

\renewcommand{\arraystretch}{1.6}
\begin{tabularx}{\linewidth}{@{} C{0.22\linewidth} C{0.18\linewidth} C{0.22\linewidth} @{}}
\rowcolor{sky}
\textbf{Time (min)} & \textbf{Frequency} & \textbf{Rel.\ Frequency} \\
\midrule
$[0, 20)$   & 4  & \ans{0.10} \\
$[20, 40)$  & 12 & \ans{0.30} \\
$[40, 60)$  & 15 & \ans{0.375} \\
$[60, 80)$  & 6  & \ans{0.15} \\
$[80, 100)$ & 3  & \ans{0.075} \\
\midrule
\rowcolor{sky}\textbf{Total} & \ans{40} & \ans{1.000} \\
\end{tabularx}

\vspace{10pt}
\begin{enumerate}[leftmargin=1.5em, itemsep=5pt]
  \item \ans{$[40, 60)$ with 15 students.}
  \item \ans{$(15 + 6 + 3)/40 = 24/40 = 0.60 = 60\%$.}
  \item \ans{Skewed right} (the tail extends toward longer reading times).
  \item \ansline{Of the 40 students surveyed, the most common reading time was
        between 40 and 60 minutes (37.5\%), and the distribution is skewed right,
        suggesting that a few students read considerably longer than typical.}
\end{enumerate}

\end{notesbox}

\vspace{0.1in}

\begin{practicebox}
\small\textbf{Choosing the Right Display --- Key}

\begin{enumerate}[leftmargin=1.5em, itemsep=8pt]
  \item \ans{Histogram.} Reason: \ans{$n = 500$ is too large for a dotplot or stemplot;
        a histogram shows the overall shape of the height distribution efficiently.}

  \item \ans{Dotplot.} Reason: \ans{$n = 8$ is very small; a dotplot shows each
        individual finishing time and makes comparisons between runners easy.}

  \item \ans{Back-to-back stemplot.} Reason: \ans{$n \approx 25$ per class is ideal
        for a stemplot; the back-to-back format lets students directly compare the
        two class distributions while preserving actual values.}
\end{enumerate}
\end{practicebox}

\vspace{0.1in}

\begin{spiralbox}
\small
\begin{multicols}{2}
\textbf{Connections forward\dots}
\begin{itemize}[itemsep=2pt]
  \item \textbf{Lesson 1.6:} SOCS framework for describing distributions.
  \item \textbf{Lesson 1.8:} Boxplots as another quantitative display.
  \item \textbf{Unit 4:} Distribution shapes $\Rightarrow$ probability models.
\end{itemize}

\columnbreak
\textbf{Big Idea:}\\
\ans{A histogram loses individual values --- you can't tell whether someone read exactly
42 or 57 minutes if both fall in $[40, 60)$. That trade-off is worth it when $n$ is
large because the shape of the distribution is more informative than any single value.}
\end{multicols}
\end{spiralbox}

\end{document}
